\chapter{Physical Design}

This chapter discusses physical optimization choices for the five required operations.
The analysis is based on \texttt{EXPLAIN PLAN} and query outputs captured in SQL Developer.
\section{Operation 1: Register new customer}

\textbf{Query type:} INSERT on CUSTOMER

\begin{lstlisting}
INSERT INTO CUSTOMER VALUES (
    CustomerTY(:taxcode, :customer_type, :first_name, :last_name, :vat, :company_name)
);
\end{lstlisting}

\begin{figure}[H]
    \centering
    \safeincludegraphics[width=0.95\textwidth]{images/op1explain.png}
    \caption{Operation 1 - EXPLAIN PLAN}
\end{figure}

\begin{figure}[H]
    \centering
    \safeincludegraphics[width=0.85\textwidth]{images/op1output.png}
    \caption{Operation 1 - Query output}
\end{figure}

\textbf{Considerations:}
\begin{itemize}
    \item This is a write operation with primary-key validation only.
    \item No additional secondary index is justified.
\end{itemize}

\section{Operation 2: Record new booking}

\textbf{Query type:} INSERT on BOOKING with \texttt{REF} resolution and trigger checks

\begin{lstlisting}
INSERT INTO BOOKING VALUES (
    BookingTY(
        :code,
        :type,
        :cost,
        :book_date,
        :interval_m,
        :is_special,
        (SELECT REF(st) FROM SETUPTEAM st WHERE st.Code = :team_code),
        (SELECT REF(el) FROM EVENTLOCATION el WHERE el.Code = :location_code)
    )
);
\end{lstlisting}

\begin{figure}[H]
    \centering
    \safeincludegraphics[width=0.95\textwidth]{images/op2explain.png}
    \caption{Operation 2 - EXPLAIN PLAN}
\end{figure}

\begin{figure}[H]
    \centering
    \safeincludegraphics[width=0.85\textwidth]{images/op2output.png}
    \caption{Operation 2 - Query output}
\end{figure}

\textbf{Considerations:}
\begin{itemize}
    \item Lookup by \texttt{SETUPTEAM.Code} and \texttt{EVENTLOCATION.Code} is already supported by primary-key indexes.
    \item Additional indexing is not needed for this insert-dominant operation.
\end{itemize}

\section{Operation 3: Register new event location}

\textbf{Query type:} INSERT on EVENTLOCATION

\begin{lstlisting}
INSERT INTO EVENTLOCATION VALUES (
    EventLocationTY(
        :code,
        :address,
        :number_h,
        :city,
        :prov,
        :zip,
        (SELECT REF(cu) FROM CUSTOMER cu WHERE cu.TaxCode = :customer_taxcode)
    )
);
\end{lstlisting}

\begin{figure}[H]
    \centering
    \safeincludegraphics[width=0.95\textwidth]{images/op3explain.png}
    \caption{Operation 3 - EXPLAIN PLAN}
\end{figure}

\begin{figure}[H]
    \centering
    \safeincludegraphics[width=0.85\textwidth]{images/op3output.png}
    \caption{Operation 3 - Query output}
\end{figure}

\textbf{Considerations:}
\begin{itemize}
    \item \texttt{CUSTOMER.TaxCode} is primary key and already indexed.
    \item No additional index is justified for this write operation.
\end{itemize}

\section{Operation 4: View team handling a specific location}

\textbf{Query type:} SELECT on BOOKING filtered by event location, ordered by date

\begin{lstlisting}
SELECT st.Code || ' - ' || st.Name
FROM BOOKING b
JOIN SETUPTEAM st ON DEREF(b.Team).Code = st.Code
WHERE DEREF(b.EventLocation).Code = :location_code
ORDER BY b.BookDate DESC
FETCH FIRST 1 ROW ONLY;
\end{lstlisting}

\begin{figure}[H]
    \centering
    \safeincludegraphics[width=0.95\textwidth]{images/op4explainbefore.png}
    \caption{Operation 4 - EXPLAIN PLAN before optimization}
\end{figure}

\begin{figure}[H]
    \centering
    \begin{subfigure}[t]{0.49\textwidth}
        \centering
        \safeincludegraphics[width=\textwidth]{images/op4outputbefore1.png}
    \end{subfigure}
    \hfill
    \begin{subfigure}[t]{0.49\textwidth}
        \centering
        \safeincludegraphics[width=\textwidth]{images/op4outputbefore2.png}
    \end{subfigure}
    \caption{Operation 4 - Query output before optimization}
\end{figure}

\subsection{Optimization}

To optimize this operation, the following index is created:

\begin{lstlisting}
CREATE INDEX idx_booking_eventlocation_date
ON BOOKING (EventLocation, BookDate DESC);
\end{lstlisting}

\begin{figure}[H]
    \centering
    \safeincludegraphics[width=0.95\textwidth]{images/op4explainafter.png}
    \caption{Operation 4 - EXPLAIN PLAN after optimization}
\end{figure}

\begin{figure}[H]
    \centering
    \begin{subfigure}[t]{0.49\textwidth}
        \centering
        \safeincludegraphics[width=\textwidth]{images/op4outputafter1.png}
    \end{subfigure}
    \hfill
    \begin{subfigure}[t]{0.49\textwidth}
        \centering
        \safeincludegraphics[width=\textwidth]{images/op4outputafter2.png}
    \end{subfigure}
    \caption{Operation 4 - Query output after optimization}
\end{figure}

\textbf{Considerations:}
\begin{itemize}
    \item Before optimization: the plan uses a full table access on BOOKING with explicit sort.
    \item After optimization: the plan switches to index access on \texttt{idx\_booking\_location\_date}.
    \item Observed plan cost decreases from about 7 to 4 in the captured execution plans.
    \item This choice is justified by the daily frequency of Operation 4.
\end{itemize}

\section{Operation 5: Rank locations by booking count}

\textbf{Query type:} aggregation on BOOKING grouped by event location

\begin{lstlisting}
SELECT
    el.Code AS LocationCode,
    COUNT(b.Code) AS NumBookings
FROM EVENTLOCATION el
LEFT JOIN BOOKING b
    ON DEREF(b.EventLocation).Code = el.Code
GROUP BY el.Code
ORDER BY NumBookings DESC, el.Code;
\end{lstlisting}

\begin{figure}[H]
    \centering
    \safeincludegraphics[width=0.95\textwidth]{images/op5explainbefore.png}
    \caption{Operation 5 - EXPLAIN PLAN before optimization}
\end{figure}

\begin{figure}[H]
    \centering
    \begin{subfigure}[t]{0.49\textwidth}
        \centering
        \safeincludegraphics[width=\textwidth]{images/op5outputbefore1.png}
    \end{subfigure}
    \hfill
    \begin{subfigure}[t]{0.49\textwidth}
        \centering
        \safeincludegraphics[width=\textwidth]{images/op5outputbefore2.png}
    \end{subfigure}
    \caption{Operation 5 - Query output before optimization}
\end{figure}

\subsection{Optimization}

No extra index is introduced beyond \texttt{idx\_booking\_eventlocation\_date}. The first key column (\texttt{EventLocation}) is already useful for grouping access paths.

\begin{figure}[H]
    \centering
    \safeincludegraphics[width=0.95\textwidth]{images/op5explainafter.png}
    \caption{Operation 5 - EXPLAIN PLAN after optimization}
\end{figure}

\begin{figure}[H]
    \centering
    \begin{subfigure}[t]{0.49\textwidth}
        \centering
        \safeincludegraphics[width=\textwidth]{images/op5outputafter1.png}
    \end{subfigure}
    \hfill
    \begin{subfigure}[t]{0.49\textwidth}
        \centering
        \safeincludegraphics[width=\textwidth]{images/op5outputafter2.png}
    \end{subfigure}
    \caption{Operation 5 - Query output after optimization}
\end{figure}

\textbf{Considerations:}
\begin{itemize}
    \item Operation 5 still requires aggregation and sorting, so a full elimination of sort cost is not expected.
    \item The same index used for Operation 4 supports the join and grouping path on booking keys.
    \item Observed plan cost decreases from about 7 to 5 in the captured execution plans.
    \item Avoiding multiple write-side indexes preserves insert/update performance for frequent operations.
\end{itemize}

\section{Index Summary}

\begin{center}
\begin{tabular}{|p{4cm}|p{5cm}|p{3cm}|}
\hline
\textbf{Index Name} & \textbf{Columns} & \textbf{Used By Operation} \\
\hline
idx\_booking\_eventlocation\_date & BOOKING(EventLocation, BookDate DESC) & Operations 4, 5 \\
\hline
\end{tabular}
\end{center}

\section{Final considerations}

The selected physical design balances read and write costs:

\begin{itemize}
    \item write-intensive operations 1, 2 and 3 remain lightweight;
    \item read operation 4 is significantly improved with a dedicated composite index;
    \item operation 5 benefits from the same index while keeping write overhead low;
    \item physical choices are coherent with the object-relational \texttt{REF/DEREF} model.
\end{itemize}



